\subsubsection{3.2 Ein ausdrückliches
Spinwort}\label{ein-ausdruxfcckliches-spinwort}

Auf dem 32-dimensionalen Hilfsfockraum seien

\[\adjustbox{max width=\linewidth}{$\displaystyle 
\gamma_j=a_j+a_j^\dagger,\qquad
\gamma_{j+5}=i(a_j^\dagger-a_j),\qquad 0\le j<5,
$}\]

und Π dessen Fermionparität. Setze

\[\adjustbox{max width=\linewidth}{$\displaystyle 
\Omega_{10}=\gamma_0\gamma_1\cdots\gamma_9=i\Pi,
\qquad
R_{ab}=\frac{(\gamma_a-\gamma_b)(\gamma_{a+5}-\gamma_{b+5})}{2}.
$}\]

R\_ab ist das Produkt zweier reeller Einheitsvektoren der
Cliffordalgebra, also ein Spin(10)-Element. Direkt gilt
\codeword{R\_ab=i Γ((ab))}. Da \codeword{p=(01)(02)(34)} mit rechts
zuerst angewandter Permutation, folgt

\[\adjustbox{max width=\linewidth}{$\displaystyle 
\boxed{S=\Omega_{10}R_{01}R_{02}R_{34}=\Pi\Gamma(p)\in\mathrm{Spin}(10).}
$}\]

S ist ein Produkt von 16 reellen Clifford-Einheitsvektoren. Die
Matrixprüfung bestätigt ohne Rundungstoleranzen

\[\adjustbox{max width=\linewidth}{$\displaystyle 
S_{even}=\Gamma(p)_{even},\qquad S_{odd}=-\Gamma(p)_{odd},
\qquad S\gamma_jS^\dagger=-\gamma_{p(j)}
$}\]

mit entsprechender Fortsetzung auf die zweite Fünfergruppe. Damit
stimmen \textbf{beide nativen Darstellungen} und die ursprüngliche
Kopplung überein:

\[\adjustbox{max width=\linewidth}{$\displaystyle 
G_F=S_{even}\otimes I_4,\qquad
G_B=(-P_{10})\otimes I_6,\qquad
W\Lambda^2G_F=G_BW.
$}\]

S hat Ordnung sechs. Die Konjugation aller 45 Spin-Erzeuger und die
Kommutation mit allen 15 Farb-Erzeugern wurden zusätzlich direkt
geprüft. Die Aussage ist stärker und präziser als ein lediglich
behaupteter Normalisatorsatz.

\textbf{Bedeutung:} Der dokumentierte Clock und die innere
Spin-Symmetrie brauchen hier keine getrennten zusätzlichen
mathematischen Träger. Es handelt sich um einen konkreten endlichen
Schritt derselben Darstellung. Daraus folgen aber weder kontinuierliche
Spin-Kontrollen noch die Identifikation des Clocks mit Hamiltonzeit. Die
Hilfs-Cliffordachsen sind nicht bereits fünf Raumrichtungen.

\subsection{4. Clock und getrennte Casimire lassen sich gemeinsam
auslesen}\label{clock-und-getrennte-casimire-lassen-sich-gemeinsam-auslesen}

Weil G aus Spin(10) stammt und auf SU(4) trivial wirkt, kommutiert es
mit beiden getrennten Casimiren, auch auf ihren Fockhebungen. Die Frage
aus dem neuen Text ist damit auf algebraischer Ebene beantwortet.

Die exakten Charakterfolgen von S\_even, S\_odd und dem Vektor sind

\[\adjustbox{max width=\linewidth}{$\displaystyle 
(16,0,4,0,4,0),\quad(16,0,4,0,4,0),\quad(10,0,4,-6,4,0).
$}\]

Die schon unabhängig nachgerechneten Zerlegungen
\codeword{Sym³16=672+144}, \codeword{Λ³16=560},
\codeword{S21(16)=1200+144+16̄} und \codeword{10⊗16=144+16̄} bestimmen die
Charaktere der sieben Typen. Die diskrete Fourierinversion erfolgt exakt
in Z{[}ζ6{]}, nicht durch gerundete komplexe Eigenwerte. Für die drei
zuvor gemeinsam dunklen Räume ergibt sich:

{\def\LTcaptype{none} % do not increment counter
\begin{longtable}[]{@{}
  >{\raggedright\arraybackslash}p{(\linewidth - 12\tabcolsep) * \real{0.1111}}
  >{\raggedleft\arraybackslash}p{(\linewidth - 12\tabcolsep) * \real{0.1481}}
  >{\raggedleft\arraybackslash}p{(\linewidth - 12\tabcolsep) * \real{0.1481}}
  >{\raggedleft\arraybackslash}p{(\linewidth - 12\tabcolsep) * \real{0.1481}}
  >{\raggedleft\arraybackslash}p{(\linewidth - 12\tabcolsep) * \real{0.1481}}
  >{\raggedleft\arraybackslash}p{(\linewidth - 12\tabcolsep) * \real{0.1481}}
  >{\raggedleft\arraybackslash}p{(\linewidth - 12\tabcolsep) * \real{0.1481}}@{}}
\toprule\noalign{}
\begin{minipage}[b]{\linewidth}\raggedright
Dunkler Typ
\end{minipage} & \begin{minipage}[b]{\linewidth}\raggedleft
Phase 0
\end{minipage} & \begin{minipage}[b]{\linewidth}\raggedleft
1
\end{minipage} & \begin{minipage}[b]{\linewidth}\raggedleft
2
\end{minipage} & \begin{minipage}[b]{\linewidth}\raggedleft
3
\end{minipage} & \begin{minipage}[b]{\linewidth}\raggedleft
4
\end{minipage} & \begin{minipage}[b]{\linewidth}\raggedleft
5
\end{minipage} \\
\midrule\noalign{}
\endhead
\bottomrule\noalign{}
\endlastfoot
(1200,20) & 4000 & 4000 & 4000 & 4000 & 4000 & 4000 \\
(560,20′) & 1920 & 1840 & 1840 & 1920 & 1840 & 1840 \\
(672,4̄) & 464 & 440 & 440 & 464 & 440 & 440 \\
Summe, bisherige Clock-Aus\-le\-sung & 6384 & 6280 & 6280 & 6384 & 6280 &
6280 \\
\end{longtable}
}

Ein separater Casimir trennt jede der sechs dunklen Clockzellen in genau
drei Teile. Die hellen Zweier-Multiplizitätsräume bleiben erhalten.
Daraus folgt im N=3-Sektor

\[\adjustbox{max width=\linewidth}{$\displaystyle 
\boxed{\dim\operatorname{Alg}^*(X,N_b,G,C_S)=96,}
\qquad
\boxed{\dim\operatorname{Alg}^*(X,N_b,G,C_S)'=119597824.}
$}\]

Zum Vergleich: ohne C\_S sind es 84 und 240742144. C\_C kann C\_S
ersetzen. Die Lagrangeprojektoren auf den drei dunklen C\_S-Werten
bestätigen die Trennung ausdrücklich; beide Casimire zugleich werden
nicht benötigt.

\textbf{Es handelt sich um eine gelöste Kombinationsfrage, nicht um eine
neue Kontrollfreigabe.} Der benötigte Casimir ist noch kein vom Compiler
bereitgestelltes Messinstrument. Der große verbleibende Kommutant zeigt
weiterhin viele nicht aufgelöste Freiheitsgrade. Die Dimensionszahlen
gelten für N=3, nicht als vollständige Klassifikation aller
Ladungssektoren.

\subsection{5. Neuer kleiner Transportzeuge mit tatsächlichen
W-Kanälen}\label{neuer-kleiner-transportzeuge-mit-tatsuxe4chlichen-w-kanuxe4len}

Ein gemeinsamer Fermion-/Boson-Dreierstern hätte die einfache Identität

\[\adjustbox{max width=\linewidth}{$\displaystyle 
Q_x=b^\dagger s f_x,\quad Q_y=b^\dagger s f_y,
\qquad [Q_y^\dagger,Q_x]=(N_b+n_s)f_y^\dagger f_x.
$}\]

Auf festem \codeword{K=Nb+n\_s≥1} ist \codeword{c†=b†s/√K} exakt eine
CAR-Mode. Das ergibt eine gewöhnliche Dreimoden-Transferkette.
\textbf{Aber dieser Stern ist nicht in einer nativen W-Zeile enthalten:}
Jede Zeile ist ein Matching aus acht disjunkten Paaren. Die Analogie
allein wäre daher eine falsche Herkunftsangabe.

Im echten W gibt es stattdessen die Paare \codeword{(4,57)} in Kanal 0
und \codeword{(4,58)} in Kanal 1. Mit den sieben besetzten
Pauli-Blockern

\[\adjustbox{max width=\linewidth}{$\displaystyle 
S_b=\{8,16,28,32,44,52,56\}
$}\]

und einem \textbf{zusätzlichen} Mischer \codeword{J(b1†b0+b0†b1)}
schließt die volle native Wechselwirkung exakt auf den vier Zuständen

\[\adjustbox{max width=\linewidth}{$\displaystyle 
S_b+\{4,57\}\ \longleftrightarrow\ S_b+b_0
\ \longleftrightarrow\ S_b+b_1
\ \longleftrightarrow\ S_b+\{4,58\}.
$}\]

Alle haben N=9. Die vollständige Wirkung sämtlicher 480 signierter
Paarterme bleibt in diesem Raum, mit Matrix

\[\adjustbox{max width=\linewidth}{$\displaystyle 
H_4=\begin{pmatrix}
0&g&0&0\\g&\Delta&J&0\\0&J&\Delta&g\\0&0&g&0
\end{pmatrix}.
$}\]

Für \codeword{g/Δ=1/20}, \codeword{J/Δ=1-1/(10√3)} und
\codeword{tΔ=20π√3} gilt analytisch

\[\adjustbox{max width=\linewidth}{$\displaystyle 
P(57\to58)>\frac{2009992727}{2022609600}
=0.9937620819\ldots>99.3\%.
$}\]

Die Eigenfrequenzen Δ±J bleiben positiv. Dies ist keine unkontrollierte
Vierzustandsabschneidung: Zuerst ist die Invarianz unter der
vollständigen Wechselwirkung bewiesen, danach wird deren exakte
Einschränkung gelöst. Ohne Mischer zerfällt sie in zwei getrennte Blöcke
und der Transfer ist null.

Der reine Einzelmischer kommutiert nicht mit dem ursprünglichen Clock.
Die Clock-Orbitsumme der Mischer \codeword{(0,1)+(12,13)+(6,7)} ist
dagegen Clock-invariant und hat auf diesem Vierzustandsraum dieselbe
Wirkung. Somit wäre „jeder solche Transfer muss den Clock brechen``
falsch. Beide Mischer sind jedoch noch zusätzliche Operationen;
Clock-Invarianz allein beweist keine Erzeugbarkeit. Ein gemeinsamer
SU(4)-Cartan Q kommutiert mit H, Nb, Clock und beiden getrennten
Casimiren. Auf den vier Zuständen besitzt er aber die Werte
\codeword{(7,7,9,9)}. Die beiden Mischer ändern diese Quantenzahl:
\codeword{||[Q\_B,M01]||²\_F=8}, \codeword{||[Q\_B,Morb]||²\_F=24}.
\textbf{Auch die Clock-invariante Orbitsumme ist damit aus diesem
Operationssatz ausgeschlossen.}

Die noch fehlende Ressource ist hier konkret: eine Operation, die diese
innere Cartanladung verändert, oder eine global hergeleitete Kopplung an
einen ausdrücklich mitgerechneten Ladungsausgleich. Mehr Wörter aus dem
unveränderten erhaltenden Alphabet können beide Mischer nicht liefern.

Dieser Zeuge ist \textbf{keine} Propagation der
N=64-Grundzustands-Lochanregung, keine räumliche Zwei-Banken-Übertragung
und keine Konstruktion ihrer Präparation. Er zeigt konstruktiv, wie nah
eine kleine Paarumwandlung an einem Transfermechanismus liegen kann, und
benennt die fehlende Ressource.

\subsection{6. Das Feldwörterbuch: eine Basisänderung allein reicht
nicht}\label{das-feldwuxf6rterbuch-eine-basisuxe4nderung-allein-reicht-nicht}

Der bestätigte Grundfehler bleibt: \codeword{M\_A⊗ε\_Lorentz} ist
symmetrisch und verschwindet als Grassmann-Bilinear. Das betrifft die
lokale ableitungsfreie Zuordnung einer gleichhändigen Weylkopie je
ursprünglichem Label bei unverändertem antisymmetrischem W. Der
nichtverschwindende gleichhändige Kanal trägt \codeword{(1,0)}; daraus
folgt noch keine gesunde Kinetik.

Der neue Fünf-Zyklus im Trägergraphen widerlegt eine rein diagonale
Zuweisung entgegengesetzter Händigkeiten an jedes Paar. Ein neuer
eigener Satz geht darüber hinaus. Für irgendeine hermitesche interne
Händigkeit Γ müsste im unveränderten Ein-Kopie-Vek\-tor\-ansatz gelten

\[\adjustbox{max width=\linewidth}{$\displaystyle 
\Gamma^TM_A+M_A\Gamma=0\quad\forall A.
$}\]

Adjungieren und Einsetzen liefert \codeword{[Γ,M\_A†M\_B]=0}. Der
ursprüngliche Tensor faktorisiert exakt als

\[\adjustbox{max width=\linewidth}{$\displaystyle 
M_{ka}=S_k\otimes C_a,\quad
\sum_kS_k^\dagger S_k=5I_{16},\quad
\sum_aC_a^\dagger C_a=3I_4.
$}\]

Die S\_k†S\_l erzeugen M16; die C\_a†C\_b erzeugen M4. Dafür wurden
vollständige Rangzertifikate modulo 101 mit 256 bzw. 16 unabhängigen
ganzzahligen Matrixwörtern konstruiert. Ein nichtverschwindender
Determinant modulo 101 beweist Nichtverschwindung über C; keine
Rangdefizienz wird übertragen. Damit ist Γ skalar, und die ursprüngliche
Gleichung erzwingt Γ=0.

\textbf{Keine hermitesche Involution Γ²=I auf demselben Träger repariert
diesen Ein-Kopie-Vek\-tor\-ansatz.} Der Ausschluss ist nun basisunabhängig,
aber weiterhin vertragsgebunden.

Nicht ausgeschlossen sind zusätzliche Dirackomponenten, ein unabhängiger
Zweierfaktor oder Ableitungsterme. Insbesondere ist

\[\adjustbox{max width=\linewidth}{$\displaystyle 
M_A\otimes\epsilon_{aux}\otimes\epsilon_{Lorentz}
$}\]

antisymmetrisch und nicht null. Ein ableitungshaltiger gleichgeladener
Vektorstrom \codeword{M\_IJ ε\_ab ψ\_Ia ↔∂\_μ ψ\_Jb} ist ebenfalls nicht
null; der pauschale Vektorausschluss durch Ladung gilt nicht außerhalb
der ableitungsfreien Klasse.

Die Zahlen 180/128 oder 60/256 zählen Komponenten eines bestimmten
Feldansatzes, nicht automatisch unabhängige propagierende Oszillatoren.
Schur fixiert interne Intertwiner, nicht alle Lorentz-/Ableitungsterme.
Der behauptete Zwang „zuerst räumlich skalieren, danach Feldtyp prüfen``
folgt nicht daraus. Zuerst bleibt ein kleiner konsistenter
Typ-/Kinetiktest sinnvoll; seine reale räumliche Herkunft muss
anschließend gemeinsam geprüft werden.

\subsection{7. Was der nachgereichte Universalraum-Text
trägt}\label{was-der-nachgereichte-universalraum-text-truxe4gt}

\subsubsection{7.1 Eine sinnvolle Hypothese, kein bereits
identifiziertes universelles
Objekt}\label{eine-sinnvolle-hypothese-kein-bereits-identifiziertes-universelles-objekt}

Nichtorthogonale oder überlappende lokale Einbettungen in eine
gemeinsame Algebra sind eine ernstzunehmende Alternative zu unabhängigen
Banken. Sie können Voraussetzungen der lokalen Paritätsschranke ändern.
Das ist ein Forschungsauftrag, kein Widerspruch zur bisherigen Schranke.

Bei Einbettungen \codeword{J\_A:V\_A→V} gilt für die CAR jedoch

\[\adjustbox{max width=\linewidth}{$\displaystyle 
\{f_A(u),f_B(v)^\dagger\}=\langle J_Au,J_Bv\rangle.
$}\]

Die Überlappungs-Grammatrix gehört also zwingend zum Modell. Die alten
Produktzustände \codeword{Ω\_A⊗Ω\_B}, ihre beiden unabhängigen
Ladungssektoren und der Zwei-Banken-Polraum dürfen nicht ungeprüft in
einen überlappenden Träger übernommen werden. Der gemeinsame
Hamiltonoperator und sein Zustand müssen dort neu beziehungsweise durch
einen bewiesenen Adapter konstruiert werden.

\subsubsection{7.2 Der vorgeschlagene Kill-Test muss korrigiert
werden}\label{der-vorgeschlagene-kill-test-muss-korrigiert-werden}

Für unabhängige orthogonale Banken kann Π\_AΠ\_B deren gemeinsame
Parität sein. Für überlappende Unterräume gilt das nicht allgemein; ihre
lokalen Paritäten müssen nicht einmal kommutieren. Der Test
\codeword{[U,Π\_AΠ\_B]=0} ist deshalb im neuen Bild nicht der richtige
universelle Test. Stattdessen muss die tatsächliche globale Parität
Π\_global separat definiert und erhalten werden. Ein
nichtverschwindender Kommutator mit einer lokalen Parität belegt zudem
allein noch keinen gerichteten Transfer.

\subsubsection{7.3 Der einfachste False-Positive: Ein Signal ohne
Bewegung}\label{der-einfachste-false-positive-ein-signal-ohne-bewegung}

Seien \codeword{a=c1}, \codeword{b=(c1+c2)/√2} auf zwei globalen
CAR-Moden und \codeword{H=ωN}. Dann haben die beiden
Einteilchen-Chartzustände bereits Überlappung \codeword{1/√2}. In ihrer
nichtorthogonalen Basis ist \codeword{K=ωS} mit einem
nichtverschwindenden Offdiagonalelement. Trotzdem bleibt die betreffende
Wahrscheinlichkeit zu jeder Zeit genau 1/2. Es wurde nichts
transportiert.

Ein belastbarer gemeinsamer Test muss daher mindestens

\[\adjustbox{max width=\linewidth}{$\displaystyle 
S_{AB}=\langle\psi_A,\psi_B\rangle,\qquad
K_{AB}=\langle\psi_A,H\psi_B\rangle
$}\]

und die tatsächliche Zeitantwort gemeinsam bestimmen. Auf einem positiv
definiten Gramträger ist
\codeword{S\textasciicircum{}\{-1/2\}KS\textasciicircum{}\{-1/2\}} der
korrekte orthonormierte Operator; bei singulärem S ist zunächst der
Nullraum zu quotientieren. Im Gegenbeispiel ergibt dies nur ωI. Die
Offdiagonale von K war kein Hop.

Das trifft sogar unmittelbar den vorhandenen nativen Pol. Dessen
64-dimensionaler Raum besitzt genau \textbf{eine} Energie E\_h, also
\codeword{P\_h H P\_h=E\_h P\_h}. Sind A und B nur neue Charts innerhalb
dieses selben Polraums, gilt zwangsläufig \codeword{K=E\_h S}. Bloßes
Umbenennen der 64 entarteten Polrichtungen erzeugt daher keine
Ausbreitung durch H. Eine tatsächliche globale Dynamik beziehungsweise
eine verfügbare aktive Operation muss hinzukommen und aus der Quelle
ausgewiesen werden.

Für die eigentliche TFPT-Lochantwort wäre der gemeinsame Gegenstand

\[\adjustbox{max width=\linewidth}{$\displaystyle 
G_{AB}(t)=\langle\Omega,
 f_B^\dagger e^{-it(H-E_0)}f_A\Omega\rangle,
$}\]

mit einer einzigen Quelle, einem H und einem Ω. \codeword{G(0)} bezahlt
die schon vorhandene Überlappung; die dynamische Änderung muss davon
unterschieden werden. Das ist die geschärfte, kleinere nächste
Forschungsaufgabe.

